\documentclass[../../main.tex]{subfiles}
\begin{document}

{\bf [解]}
題意より
\begin{align}
 (x^3+ax^2+bx+c)^2=(x^2-1)(x^2+px+q)^2+D \label{eq:1}
\end{align}
が $x$ についての恒等式であるような$D$を求めれば良い．
\cref{eq:1}の左辺，右辺を各々 $f(x),g(x)$ とすると
\begin{align}
 \begin{cases}
 f(x)=x^6+2ax^5+(2b+a^2)x^4+(2c+2ab)x^3+(2ac+b^2)x^2+2bcx+c^2 \\
 g(x)=x^6+2px^5+(p^2+2q-1)x^4+(2pq-2p)x^3+(q^2-p^2-2q)x^2-2pqx-q^2+D  
 \end{cases}
\end{align}
だから係数比較して，
\begin{align}
  & a=p  \label{eq:2} \\
  & 2b=2q-1 \label{eq:3} \\
  & ab+c=pq-p \label{eq:4} \\
  & 2ac+b^2=q^2-p^2-2q \label{eq:5} \\
  & bc=-pq \label{eq:6} \\
  & c^2=-q^2+D \label{eq:7}
\end{align}
を得る．
したがって\cref{eq:2,eq:3}より
\begin{align}
a=p, \qquad b=q-\dfrac{1}{2} \label{eq:8}
\end{align}
である．これを\cref{eq:4}に代入して
\begin{align}
c=-\dfrac{1}{2} p \label{eq:9}
\end{align}
だから，\cref{eq:8,eq:9}を\cref{eq:5}に代入して
\begin{align}
 q=-\dfrac{1}{4} \label{eq:10}
\end{align}
となり，これを\cref{eq:8}に代入して
\begin{align}
 b=-\dfrac{3}{4} \label{eq:11}
\end{align}
となる．\cref{eq:6}から
\begin{align}
bc=-pq \quad\Longrightarrow\quad \dfrac{3}{8}p=\dfrac{1}{4}p \quad\therefore\ p=0
\end{align}
だから，\cref{eq:8,eq:9}より $a=c=0$ で\cref{eq:7}に代入して
\begin{align}
D=c^2+q^2=0+\dfrac{1}{16}=\dfrac{1}{16}
\end{align}
が求める値である．


\newpage

\end{document}
